\section{Tanimoto dot product kernel and its random features}\label{sec:t-dp}

\citet{ralaivola2005graph} gave a definition for the Tanimoto coefficient involving dot products:
\begin{equation}\label{eqn:tanidot-kernel}
    T_{DP}(x, x') = \frac{x\cdot x'}{\|x\|^2+\|x'\|^2-x\cdot x'},
\end{equation}
with $T_{DP}(x,x')=1$ when $x,x'=0$. It is easy to check that $T_{DP}(x,x')=T_{MM}(x,x')$ on binary vectors, which was used by \citet{ralaivola2005graph} to prove that $T_{DP}$ is a kernel on the space $\{0,1\}^d$, referencing prior work by \citet{gower1971general}.  However, $T_{DP}$ is not identical to $T_{MM}$ for general inputs $x,x'\in \R_{\ge 0}^d$. Here, we give the first proof that $T_{DP}$ is a positive definite function in $\R^d$ and thus, also a valid kernel in this space.

\begin{theorem}\label{thm:tdp-is-pd}
    For $x,x'\neq 0$ in $\R^d$, we have
    \begin{equation}\label{eqn:tanidot-power-series}
       T_{DP}(x,x')= \sum_{r=1}^\infty \left(x\cdot x'\right)^r \left(\|x\|^2+\|x'\|^2\right)^{-r},
    \end{equation}
    where the series is absolutely convergent. The function $T_{DP}$ is a positive definite kernel in $\R^d$.
\end{theorem}

It has been noticed previously that, unlike $1-T_{MM}$, the function $1-T_{DP}$ is \emph{not} a distance metric on non-binary inputs \citep{kosub2019note}. Indeed, when $d=1$, the inputs $\{1,2,4\}$   violate the triangle inequality. However, we can easily derive a distance metric from $T_{DP}$.

\begin{corollary}\label{thm:tdp-is-metric}
    $d_{DP}(x,x')=\sqrt{1-T_{DP}(x,x')}$ 
    corresponds to the RKHS norm of the function $\frac{1}{2}[T_{DP}(x,\cdot)-T_{DP}(x',\cdot)]$
    and is therefore a valid distance metric on $\mathbb{R}^d$.
\end{corollary}

Proofs are given in Appendix~\ref{apdx:tdp-is-pd-proof}. These results imply that $T_{DP}$, like $T_{MM}$, is an extension of the set-valued Tanimoto coefficient (equation~\ref{eqn:tanimoto-for-sets})
to real vectors and can be used as a substitute for $T_{MM}$ in machine learning algorithms that require a kernel or distance metric. Unlike $T_{MM}$, the kernel $T_{DP}$
is differentiable everywhere with respect to its inputs.
It can also be computed in batches using matrix-matrix multiplication,
allowing for efficient vectorized computation.

We now consider producing a random features approximation to $T_{DP}$ for large-scale applications.
Motivated by the close relationship between $T_{DP}$ and $T_{MM}$,
one may be tempted to find a random hash for $T_{DP}$
and apply the techniques developed in section~\ref{sec:minmax-rf}.
Unfortunately, we are able to prove that this is not possible.
\begin{proposition}\label{no random hash for tdp}
    There exists no random hash function for $T_{DP}$ over non-binary vectors.
\end{proposition}
\begin{proof}
    \citet{charikar2002similarity} proved that if $s(x,x')$ is a similarity function
    for which there exists a random hash,
    then $1-s(x,x')$ must satisfy the triangle inequality (see their Lemma 1).
    Because $1-T_{DP}(x,x')$ does not satisfy the triangle inequality it follows by contradiction
    that there does not exist a random hash for $T_{DP}$.
\end{proof}

Therefore producing random features for $T_{DP}$ will require another approach.
In the remainder of this section we present a framework to produce random features for $T_{DP}$ by directly approximating its power series
(equation~\ref{eqn:tanidot-power-series}).
We first describe a method to produce random features for $(\|x\|^2+\|x'\|^2)^{-r}$ (\ref{ssec:prefactor-rf}).
Then we describe how these features can be combined with existing random features for the polynomial kernel
to approximate $T_{DP}$'s truncated power series (\ref{sec:tdp-rf}--\ref{ssec:implementing tdp rf}).
Lastly, we present an error bound for the kernel matrix of a dataset in the spectral norm,
showing that the required dimension for the sketch scales optimally with the stable rank of the kernel matrix
(\ref{sec:tdp-error-bound}).


\subsection{Random features for the ``prefactor'' $\left(\|x\|^2+\|x'\|^2\right)^{-r}$}\label{ssec:prefactor-rf}

In this section we present a random feature map for the positive definite kernel $\textstyle (x,x')\mapsto (\|x\|^2 + \|x'\|^2)^{-r}$,
which we will refer to as the the \emph{prefactor}.
We defer all proofs to Appendix~\ref{apdx:prefactor-sketch}.
We begin with the following lemma, which defines scalar random features for the prefactor:

\begin{lemma}\label{scalar prefactor features lemma}
    If  $Z\sim\mathrm{Gamma}(s,c)$ (where $c$ is a \emph{rate} parameter), then
    \numeq{
        \varphi_{r,Z}(x) = e^{(1/2-\|x\|^2)Z} Z^{(r-s)/2}\sqrt{c^{-s}e^{(c-1)Z}\Gamma(s)/{\Gamma(r)}}
    } is an unbiased scalar random feature for the prefactor $(\|x\|^2 + \|x'\|^2)^{-r}$
    for all $s,c>0$.
\end{lemma}
Although independent copies of $Z$ could be combined to form an $M$-dimensional sketch,
we instead propose to use a dependent point set $Z_1,\dots,Z_M$ where each element $Z_i$ has a Gamma$(s,c)$ distribution whilst maximally covering the real line. This is a well-established Quasi-Monte Carlo (QMC) technique which generally attains
lower variance. 
We define our $M$-dimensional QMC features in the following lemma:
\begin{lemma}\label{lemma:def QMC prefactor features}
     Let $\gamma_{s,c}$ be the inverse cumulative distribution function of a $\mathrm{Gamma}(s,c)$ random variable.
     Fix $M,r\in\mathbb{N}$, $u\in(0,1)$, $c,s>0$ and let $u_i = u + i/M - \lfloor u + i/M\rfloor$ for $i=1,\dots,M$.
     Define $\phi_{u,r}(x) = (\phi_{u,r,1}(x),\dots,\phi_{u,r,M}(x))$, where:
     \begin{equation}
         \phi_{u,r,i}(x) =\frac{1}{\sqrt M}\sqrt{\frac{c^{-s}\Gamma(s)}{\Gamma(r)} } e^{-(\|x\|^2-c/2)\gamma_{s,c}(u_i) }(\gamma_{s,c}(u_i))^{(r-s)/2}\ .
     \end{equation}
     If $u\sim\mathcal U(0,1)$ then $\phi_{u,r}(x)$ forms unbiased random features of the prefactor $(\|x\|^2 + \|x'\|^2)^{-r}$.
\end{lemma}
Although the random features are unbiased for all $s,c>0$, the value of these parameters will impact the error.
We show that if $s,c$ are suitably tuned, then the relative error can be bounded:
 \begin{lemma} \label{QMC error}
Let $x^{(1)},\dots,x^{(n)}\in \R^d$ with $\frac{\min_i \|x^{(i)}\|^2}{\max_i \|x^{(i)}\|^2} \ge \zeta$, and fix $u\in[0,1]$. Define the relative error
\numeq{ \label{error matrix}
E_{i,j} = \frac{   \phi_{u,r}(  x^{(i)} )\cdot \phi_{u,r}( x^{(j)} ) - (\| x^{(i)} \| ^2+ \| x^{(j)} \| ^2)^{-r} }{( \| x^{(i)} \| ^2+ \| x^{(j)} \| ^2)^{-r}} .
}
If $c = 2\zeta^2$, $s=r\zeta$, then 
for some constant $C$ independent of $r$
this error satisfies, 
\numeq{
\max_{1\le i,j \le n} | E_{i,j}| \leq \frac{2}{M} \frac{\Gamma(r\zeta)\zeta^{-r\zeta}}{\Gamma(r)} (r/e)^{r(\zeta-1)} (1.3)^r \leq C (M\zeta)^{-1}.
}
\end{lemma}

Together, these lemmas suggest random features for the prefactor can be created
by first estimating $\zeta$ (the minimum ratio of norms of input vectors),
then using the random features from \Cref{lemma:def QMC prefactor features} with the
values of $s,c$ specified in \Cref{QMC error}.












\subsection{A framework to produce random features for $T_{DP}$}
\label{sec:tdp-rf}

There are straightforward rules for producing random features for sums and products of kernels whose individual random features are known \citep[sec.\@ 2.6.2]{duvenaud_2014}.
Random features for kernels $k_1,k_2$ can be \emph{concatenated} (denoted $\oplus$) to form random features for the sum kernel $k_1+k_2$,
while their \emph{tensor product}\footnote{
For two vectors $x\in\R^{d_1}$ and $y\in\R^{d_2}$, define the tensor product $x\otimes y = \text{vec}(xy^T) \in \R^{d_1d_2}$.
} (denoted $\otimes$)
forms random features for the product kernel $k_1\times k_2$.
Our strategy to produce features for $T_{DP}$ is
to combine random features for the ``prefactor'' (presented in \cref{ssec:prefactor-rf})
with random features for the polynomial kernel to produce random features for $T_{DP}$'s power series (equation~\ref{eqn:tanidot-power-series}) truncated at $R$ terms.

Fix $R\in\mathbb{N}$, and for $r=1,\ldots,R$, let $\phi_r$ be a $m_r$-dimensional random features map for the prefactor
$\textstyle (\|x\|^2 + \|x'\|^2)^{-r}$ and let $\psi_r$ be a $m'_r$-dimensional random features map for
$\left(x\cdot x'\right)^r$. The function:
\begin{equation}\label{naive tensor product rfs}
    \tilde\Phi_R(x) = \oplus_{r=1}^R \left[\phi_r(x) \otimes \psi_r(x)\right]
\end{equation}
is therefore a random feature estimate for $T_{DP}$'s power series, truncated at $R$ terms.
Unfortunately, these random features have dimension 
$M=\sum_{r=1}^R m_r m'_r$
which depends on the \emph{product} of the random features dimension of $\phi_r$ and $\psi_r$. Furthermore, the dimension $m_r'$ of the random features $\psi_r$ required to approximate the polynomial kernel $(x\cdot x')^r$ with good accuracy can scale poorly with $r$.
For even modest values of $m_r,m'_r$ the resulting value of $M$ will likely be prohibitively large.

To remedy this, we turn to recent works which propose powerful linear maps
to approximate tensor products
with a lower-dimensional vector.
Assuming $x^{(1)},y^{(1)}\in\R^{d_1}$ and $x^{(2)},y^{(2)}\in\R^{d_2}$,
these maps are effectively random matrices $\Pi\in \R^{m\times (d_1d_2)}$,
which exhibit a \emph{subspace embedding property} whereby $[\Pi(x^{(1)}\otimes x^{(2)})]\cdot[\Pi(y^{(1)}\otimes y^{(2)})]$ concentrates sharply around $(x^{(1)}\otimes x^{(2)})\cdot(y^{(1)}\otimes y^{(2)})$.
Critically, the product $\Pi(x^{(1)}\otimes x^{(2)})$ can be computed \emph{without} instantiating either matrix $\Pi$ or the tensor product $x^{(1)}\otimes x^{(2)}$.
Examples of such methods include
\textsc{TensorSketch} and \textsc{TensorSRHT}
\citep{pagh2013compressed,pham2013fast, ahle2020oblivious},
but for generality we will simply refer to these methods as \textsc{Sketch}.
Defining a series of sketches $\textsc{Sketch}_r:\R^{m_r}\times \R^{m'_r}\mapsto \R^{\tilde{m}_r}$
we can modify $\tilde\Phi_R$ from equation~\ref{naive tensor product rfs} into:
\begin{equation}\label{sketch tensor product rfs}
    \Phi_R(x) = \oplus_{r=1}^R \textsc{Sketch}_r\left[\phi_r(x),  \psi_r(x)\right]
\end{equation}
which has output dimension $M=\sum_{r=1}^R \tilde{m}_r$, i.e.\@ without any pathological dependencies on the dimensions of $\phi_r(x),\psi_r(x)$.
However, because these features approximate a \emph{truncated} power series, they will be biased downward 
due to the monotonicity of the power series (\ref{eqn:tanidot-power-series}).
We propose two bias correction techniques to potentially improve empirical accuracy.
One approach is to normalize the random features such that $\Phi(x)\cdot\Phi(x)=1 = T_{DP}(x,x)$ for all $x\in\R^d$, i.e., such that the diagonal entries of the kernel matrix $K$ are estimated exactly.
A second approach is based on sketching the residual of the power series.
These approaches are presented in detail in Appendix~\ref{apdx:bias-correction}.


\subsection{Implementing the random features}\label{ssec:implementing tdp rf}

Instantiating the random features from the previous subsection (\ref{sketch tensor product rfs})
requires making concrete choices for $R$, $\textsc{Sketch}_r,m_r,m'_r,\tilde{m}_r,\phi_r,\psi_r$ for all $r$,
and choosing a bias correction technique.
There are many reasonable choices for $\textsc{Sketch}_r$, such as
\textsc{TensorSketch} \citep{pham2013fast}
and \textsc{TensorSRHT} \citep{ahle2020oblivious}.
These sketches generally allow $m_r,m'_r,\tilde{m}_r$ to be chosen freely
(although naturally error will increase as $\tilde{m}_r$ decreases).
Many of these sketches can also be used as random features for the polynomial kernel $\psi_r$ \citep{wacker2022improved},
either directly or as part of more complex algorithms like  \textsc{TreeSketch} \citep{ahle2020oblivious}
or complex-to-real sketches \citep{wacker2023complextoreal}.
The QMC random features from section~\ref{ssec:prefactor-rf} can be used for the prefactor $\phi_r$,
with the parameters $s,c$ chosen based on the anticipated norms of the input vectors.

The only remaining inputs are $R$ (the number of power series terms to approximate)
and $\tilde{m}_1,\ldots,\tilde{m}_R$ (how many random features to use for each term).
Assuming a fixed dimension $M$ for the final random features,
this choice involves a bias variance trade-off,
as a higher value of $R$ will reduce bias but require each term in the power series
to have fewer features, thereby increasing variance.
Intuitively, because the terms of the power series decrease monotonically
the variance of terms for small $r$ is likely to dominate the overall variance,
and therefore we surmise that a \emph{decreasing} sequence for $\{\tilde{m}_r\}_{r=1}^R$
will be the best choice.
Ultimately however we do not have theoretical results to dictate this choice in practice.
We will evaluate these choices empirically in section~\ref{sec:experiments}.











\subsection{Asymptotic error bound for $T_{DP}$ random features}\label{sec:tdp-error-bound}

Because many kernel methods use kernel matrices as linear operators,
it is natural to examine the error of kernel approximations in the operator norm.
Previous works have produced asymptotic error bounds
of the random feature dimension $m$ required to achieve a
relative approximation error of $\varepsilon$ in the operator norm.
Defining $\tilde{\text{sr}}(K) = \Tr(K)/\|K\|_\text{op}$,
\citet{cohen2015optimal} show that $m=\tilde\Omega(\tilde{\text{sr}}(K)/\varepsilon^2)$
is essentially optimal for data-oblivious random features of linear kernels, even though it is possible to eliminate logarithmic factors.
Our main theoretical result is that with the correct choices of base random features,
the random features for $T_{DP}$ presented in section~\ref{sec:tdp-rf} achieve similar scaling.
We now state this as a theorem.

\begin{theorem} \label{error bound theorem}
For any $n\ge 1$, let $x^{(1)},\dots,x^{(n)}\in\R^d$ be a set of inputs with $\frac{\min_i \|x^{(i)}\|^2}{\max_i \|x^{(i)}\|^2} \geq \zeta$. Let $K$ be the matrix with entries $K_{i,j} = T_\text{DP}(x^{(i)},x^{(j)})$. For all $\varepsilon>0$, there exists an oblivious sketch $\Phi:\R^d \to \R^m$ with $m=\tilde\Omega(\tilde{\text{sr}}(K) /\varepsilon^{2})$, such that 
\begin{equation}
\Pr_{\Phi}\big(\|\widehat K-K\|_\text{op} \ge\varepsilon \|K\|_\text{op}\big) \leq \frac{1}{\text{poly}(n)}
\end{equation}
where $\widehat K_{i,j} = \Phi(x^{(i)})\cdot \Phi(x^{(j)})$. Furthermore, the sketch can be computed in time $\tilde O(\tilde{\text{sr}}(K)n\varepsilon^{-2}+\text{nnz}(X)\varepsilon^{-2}+n\zeta^{-1}\varepsilon^{-3})$.  %
\end{theorem}
The random features in the theorem follow equation~\ref{sketch tensor product rfs},
with specific choices for $M$,$R$, and $\{\textsc{Sketch}_r,\tilde{m}_r,\phi_r,\psi_r\}_{r=1}^R$
given in the proof in Appendix~\ref{apdx:proof-of-error-bound}.
\Cref{error bound theorem} essentially suggests that, with the correct settings,
the error of the random features proposed in section~\ref{sec:tdp-rf}
scales as well as one could reasonably expect for a kernel of this type.
We would highlight that the computational cost of the sketch is sub-quadratic in $n$, and compares favourably with the cost of data-dependent low-rank approximation methods.
